% Prefácio
\chapter{Prefácio}

\lettrine{O}{ estudo da matemática} não se limita apenas às aplicações teóricas dentro de seu próprio conteúdo, ela também serve de ferramenta para muitas outras áreas do conhecimento, entre as quais: química, física, biologia, geociências, engenharias, computação, arquitetura, história etc. Podemos fazer um paralelo do nosso desenvolvimento dentro das áreas do conhecimento com a construção das pirâmides do Egito ou, até mesmo, da construção da sua casa. O primeiro passo é construir uma base/fundação (fundamental) forte e consolidada para, só então, adicionarmos novas estruturas e elementos apoiados por cima.  

Façamos uma analogia com o conhecimento: no ensino básico aprendemos o alfabeto, os números, a ler e a escrever. A partir desse momento, começamos a consolidar o que aprendemos em casa desde que nascemos, formando a base da nossa pirâmide. Logo em seguida, no ensino fundamental, começamos a nos aprofundar mais em cada conteúdo, novas áreas vão surgindo e, assim, a dificuldade e a quantidade de assuntos que temos que aprender também aumenta. Porém, temos um foco, que é nos preparar para o ensino médio: nossa pirâmide se estreita um pouco e subimos mais um andar. 

Chegando no temido ensino médio, temos diversas áreas do conhecimento para estudar em um nível relativamente aprofundado, daí surgem muitas dificuldades devido à alta quantidade de conteúdo que temos que absorver, mas, novamente, estamos focados em um objetivo: vestibulares/ENEM.

Um vez concluída essa fase, temos a graduação (faculdade/universidade), depois temos a pós graduação, mestrados, doutorados, pós-doutorados e especializações. Perceba que cada vez estamos estudando assuntos mais específicos e, portanto, afinando a nossa pirâmide até chegar ao topo (entenda como o topo da nossa pirâmide metafórica de conhecimentos específicos, e não como um ponto de prestígio social). Muitas vezes não podemos concluir ou passar por todas as fases desse processo mas, mesmo assim, podemos usar essa analogia da nossa pirâmide. Quando começamos a trabalhar, a tendência é que nos tornemos melhores e mais eficientes naquilo que fazemos, que também serve como os níveis na nossa pirâmide.

Quando a fundação/base de uma casa (ou da nossa pirâmide) não é bem feita, tudo que está construído por cima se torna frágil e corre o risco de desmoronar. Analogamente, na grande maioria das vezes, quando temos dificuldade em uma matéria/conteúdo mais específico, o nosso verdadeiro problema na verdade está na base que não foi consolidada apropriadamente. 

Deixando de lado as analogias, podemos perceber que a nossa dificuldade em resolver um exercício de física, por exemplo, pode estar na verdade mais relacionada à interpretação do problema do que com a matéria propriamente, ou mesmo em como desenvolver a parte matemática da questão (por não saber as operações e funções que são necessárias). Isso é mais frequente no ensino médio e pré-vestibular, nos quais já iniciamos com muitos conhecimentos não consolidados (do ensino fundamental) e uma carga grande de matérias para estudarmos. Estamos focados em um objetivo determinado e, por muitas vezes, deixamos de aprender/entender e apenas decoramos grande parte dos conteúdos para utilizá-los durante a prova, o que não é muito proveitoso a médio e longo prazo.

Portanto, a existência desta apostila se justifica como uma tentativa de revisar conteúdos fundamentais da matemática e ajudá-los(as) a absorver melhor alguns dos conceitos que foram apenas decorados e pouco praticados ao longo da vida, porém extremamente necessários e importantes para exercícios e conteúdos mais ``complexos''. 

\vspace{8cm}

\textbf{Caro(a) aluno(a), todos nós possuímos facilidades e dificuldades de aprendizado individuais e percepções diferentes em cada assunto. Como o objetivo dessa apostila é revisar vários dos temas fundamentais da Matemática, você pode ocasionalmente se sentir mais seguro(a) e confiante em pular determinado tema, o que não é problema algum. Porém, recomendamos que seja realizada ao menos uma leitura global do capítulo antes de ignorá-lo, pois nele pode existir um novo conceito ou técnica ainda desconhecidos por você e que podem ser muito úteis no futuro.}

\cleardoublepage   % Make sure contents page starts on right-side page